\chapter{Regulators and Licences}\label{ch:fm:regulators-and-licences}

A firm that trades only its own money need not be a client of anyone, and it has no clients of its own. Whether it needs a licence is another
matter. The second Markets in Financial Instruments Directive keeps an exemption for firms dealing on their own account in the European Union,
and then takes it away from almost every firm this series describes: from market makers, from members of a trading venue and firms with direct
electronic access to one, and from any firm that applies a high-frequency algorithmic trading technique, however small it is and although it has
no clients. Where a firm may trade, and what it must be to do so, is set by rules that differ by country and by activity, and the map of
them is part of the firm's design.

\section{Why a trading firm needs permission}

Regulation follows activities, not firms. The same firm can be unregulated for one thing it does, registered for another and licensed for a third,
and a new desk, product or office can move it across a line it did not know was there.

\begin{definition}[Regulatory perimeter, authorisation]\label{def:fm:regulators-and-licences:perimeter}
The \emph{regulatory perimeter}\index{regulatory perimeter} of a jurisdiction is the boundary between the activities that require permission from
a regulator and those that do not, drawn by the legislation's definitions of regulated activities and its exemptions. \emph{Authorisation}\index{authorisation}
is the permission, granted by a regulator on application, to carry on specified regulated activities, with the conditions, capital and
supervision that go with it.
\end{definition}

\begin{definition}[Dealing on own account, investment firm]\label{def:fm:regulators-and-licences:own}
\emph{Dealing on own account}\index{dealing on own account} is trading against proprietary capital that results in transactions in financial
instruments. An \emph{investment firm}\index{investment firm} is, in European law, a legal person whose regular occupation or business is providing
investment services to third parties or performing investment activities on a professional basis; dealing on own account is one of those
activities.
\end{definition}

The chapter's working tool is a map of these rules as data, each row naming a jurisdiction, an activity, the regime it touches, the legal
reference, the statutory decision period where there is one, the month it was checked and the ledger rows that source it
(\cref{lst:fm:regulators-and-licences:map}). The map is information about where the rules are; whether a particular firm is inside a perimeter is a
question for its lawyers and its regulator.

\omcode{code/firm/regmap/firm_regmap.py}{39}{64}{The regulatory map's schema check, and the query of a firm's profile of activities and
jurisdictions.}\label{lst:fm:regulators-and-licences:map}

\section{The United States: the SEC, the CFTC and the self-regulatory organisations}

The United States divides the field by instrument. Securities fall to the Securities and Exchange Commission, futures, options on futures and swaps
to the Commodity Futures Trading Commission; each delegates much of the registration and day-to-day supervision to a membership body.

\begin{definition}[Self-regulatory organisation]\label{def:fm:regulators-and-licences:sro}
A \emph{self-regulatory organisation}\index{self-regulatory organisation} is a membership body, such as a securities association or an exchange,
that writes and enforces rules for its members under the supervision of a government regulator, and to which registered firms are required to
belong.
\end{definition}

\begin{definition}[Swap dealer, commodity pool operator, introducing broker]\label{def:fm:regulators-and-licences:cftc}
A \emph{swap dealer}\index{swap dealer} holds itself out as a dealer in swaps, makes a market in them, or regularly enters into swaps with
counterparties for its own account in the ordinary course of business. A \emph{commodity pool operator}\index{commodity pool operator} operates a
pool that combines funds from several persons to trade futures, options on futures or swaps, and solicits funds for it. An \emph{introducing
broker}\index{introducing broker} solicits or accepts orders in futures, options on futures or swaps but does not accept money or other assets
from customers to support them.
\end{definition}

\begin{dated}{2026-09}{US registration}\label{dat:fm:regulators-and-licences:us}
\textbf{Securities.} A broker or dealer must register with the SEC to effect securities transactions (Exchange Act s. 15(a)); within 45 days of
an application the SEC grants registration or opens proceedings, which must end within 120 days, extendable by up to 90 (s. 15(b)(1)). A
registered broker-dealer must also join a registered securities association unless it trades solely on exchanges of which it is a member (s.
15(b)(8)); Rule 15b9-1, as amended in September 2023, limits that exemption to exchange members with no customer accounts that trade solely on
their exchanges, with narrow exceptions. An investment adviser must register under the Advisers Act (s. 203(a)), on the same 45-day timetable.
\textbf{Derivatives.} The Commodity Exchange Act requires certain firms to register with the CFTC, which has delegated registration to the National
Futures Association; with few exceptions registered firms must be NFA members. The categories include futures commission merchants, \omterm{def:fm:regulators-and-licences:cftc}{introducing
brokers}, \omterm{def:fm:regulators-and-licences:cftc}{commodity pool operators}, commodity trading advisors and \omterm{def:fm:regulators-and-licences:cftc}{swap dealers}.
\end{dated}

\begin{omfigure}
\begin{tikzpicture}[font=\footnotesize,>=stealth,box/.style={draw,rounded corners,align=center,minimum height=0.75cm,text width=3.2cm}]
\node[box,fill=gray!15] (sec) at (-3.4,1.6) {SEC\\{\scriptsize securities}};
\node[box,fill=gray!15] (cftc) at (3.4,1.6) {CFTC\\{\scriptsize futures, options on futures, swaps}};
\node[box,fill=omDef!15] (finra) at (-3.4,0) {securities association\\{\scriptsize membership rules}};
\node[box,fill=omDef!15] (nfa) at (3.4,0) {NFA\\{\scriptsize registration delegated}};
\node[box,fill=omProp!15] (bd) at (-3.4,-1.6) {broker-dealer, investment adviser};
\node[box,fill=omProp!15] (fcm) at (3.4,-1.6) {FCM, IB, CPO, CTA, swap dealer};
\draw[->] (sec) -- node[left,font=\scriptsize] {oversees} (finra); \draw[->] (cftc) -- node[right,font=\scriptsize] {oversees} (nfa);
\draw[->] (sec.south west) to[bend right=30] node[left,font=\scriptsize] {registers} (bd.north west);
\draw[->] (finra) -- node[left,font=\scriptsize] {members} (bd); \draw[->] (nfa) -- node[right,font=\scriptsize] {registers, members} (fcm);
\end{tikzpicture}

\omcaption{The US division of the field: the SEC for securities and the CFTC for derivatives, each with a membership body to which registered firms
must belong. Schematic of the rules in the dated box.}\label{fig:fm:regulators-and-licences:us}
\end{omfigure}

For a proprietary trading firm the 2023 amendment is the hook in American form: a firm that trades off-exchange, even for itself, now needs the
securities association's membership and its rulebook. Market access (Rule 15c3-5, Book 11) adds pre-trade controls whoever sponsors the firm's
orders.

\section{Europe and the United Kingdom}

\begin{definition}[Passporting]\label{def:fm:regulators-and-licences:passport}
\emph{Passporting}\index{passporting} is the right of a firm authorised in one member state of the European Union or the European Economic Area to
provide the services its \omterm{def:fm:regulators-and-licences:perimeter}{authorisation} covers in the others, on notification to its home regulator, without a further \omterm{def:fm:regulators-and-licences:perimeter}{authorisation}.
\end{definition}

\begin{dated}{2026-09}{EU authorisation under MiFID II}\label{dat:fm:regulators-and-licences:eu}
Directive 2014/65/EU requires prior \omterm{def:fm:regulators-and-licences:perimeter}{authorisation} for investment services and activities performed as a regular occupation (art. 5), and the
regulator must decide within six months of a complete application (art. 7(3)); an authorised firm may provide its services in other member states
(art. 34). Firms dealing only on own account are exempt (art. 2(1)(d)) unless they are market makers, members or participants of a regulated
market or MTF, have direct electronic access to a trading venue, apply a high-frequency algorithmic trading technique, or deal on own account
when executing client orders. A firm engaged in algorithmic trading must have effective systems and risk controls, business continuity and tested
systems (art. 17(1)), and must notify the authorities of its home state and of the venues where it trades (art. 17(2)).
\end{dated}

\begin{dated}{2026-09}{UK permission}\label{dat:fm:regulators-and-licences:uk}
Under the Financial Services and Markets Act 2000 no person may carry on a regulated activity in the United Kingdom unless authorised or exempt
(s. 19, the general prohibition). An application for permission must be determined within six months of the regulator receiving it complete, and
an incomplete one within twelve months (s. 55V). The UK's exemptions for own-account dealing are its own and are not restated here: the map marks
the UK rows as checks.
\end{dated}

The European exemption is worded so that almost no electronic trading firm meets it. A proprietary firm that is a member of an exchange, or that
reaches one through direct electronic access, is out; one that co-locates, generates its orders without human intervention and sends many
messages is applying a high-frequency technique (Book 10's definition in the directive's words, article 4(1)(40)) and is out again. What it gains
in return is the passport.

\begin{omfigure}
\begin{tikzpicture}[font=\footnotesize,>=stealth]
\node[draw,rounded corners,fill=gray!10,align=center,text width=4.6cm,minimum height=0.8cm] (q0) at (0,0) {deals only on its own account?};
\node[draw,rounded corners,fill=omDef!15,align=center,text width=4.6cm,minimum height=0.8cm] (q1) at (0,-1.9) {market maker, venue member, direct electronic access, HFT technique or client orders?};
\node[draw,rounded corners,fill=omThm!20,align=center,text width=2.6cm,minimum height=0.8cm] (ex) at (5.0,-1.9) {exempt (art. 2(1)(d))};
\node[draw,rounded corners,fill=omProp!20,align=center,text width=4.6cm,minimum height=0.8cm] (au) at (0,-3.9) {authorisation as an investment firm (arts. 5, 7): decision within six months};
\node[draw,rounded corners,fill=omProp!10,align=center,text width=2.6cm,minimum height=0.8cm] (pp) at (5.0,-3.9) {passport to other member states (art. 34)};
\node[draw,rounded corners,fill=omProp!20,align=center,text width=2.6cm,minimum height=0.8cm] (au2) at (-5.0,-1.9) {authorisation for its other services};
\draw[->] (q0) -- node[right,font=\scriptsize] {yes} (q1); \draw[->] (q0) -| node[above,pos=0.2,font=\scriptsize] {no} (au2);
\draw[->] (q1) -- node[above,font=\scriptsize] {no} (ex); \draw[->] (q1) -- node[right,font=\scriptsize] {yes} (au); \draw[->] (au) -- (pp);
\end{tikzpicture}

\omcaption{The own-account exemption of MiFID II for a firm trading financial instruments other than commodity derivatives, as the directive
words it. Almost every electronic trading firm answers yes to the second question. Schematic of the legal text, not advice.}\label{fig:fm:regulators-and-licences:tree}
\end{omfigure}

\section{Asia-Pacific}

Asia's markets are regulated country by country. In Hong Kong a corporation carrying on a business in a regulated activity needs a licence from
the Securities and Futures Commission; the Securities and Futures Ordinance lists thirteen types, among them dealing in securities (type 1),
dealing in futures contracts (type 2), providing automated trading services (type 7) and asset management (type 9). Whether a firm trading only its
own money is carrying on such a business turns on the ordinance's definitions and exemptions, so the map marks it as a check, not a requirement.
Singapore, Japan and Australia each have their own licences; this chapter's map carries only the rows it has sourced, and says so.

\section{Rules for algorithmic trading}

Once inside the perimeter, an electronic firm carries obligations specific to algorithms. They are the controls of Book 11's chapter on risk
controls, written into law: in the European Union, article 17 of MiFID II and its technical standards; in the United States, the market access
rule for the broker whose connection the firm uses (\cref{tab:fm:regulators-and-licences:algo}).

\begin{table}[ht]
\centering\small
\begin{tabular}{llp{6.2cm}}
\toprule
where & obligation & reference \\
\midrule
EU & systems and risk controls, limits, no erroneous orders & Directive 2014/65/EU art. 17(1) \\
EU & business continuity; tested and monitored systems & Directive 2014/65/EU art. 17(1) \\
EU & notify home and venue authorities & Directive 2014/65/EU art. 17(2) \\
EU & detailed organisational requirements & Delegated Regulation (EU) 2017/589 (Book 11) \\
US & pre-trade risk controls for market access & Rule 15c3-5 (Book 11) \\
\bottomrule
\end{tabular}
\caption{The algorithmic-trading checklist of \texttt{firm.regmap.ALGO\_CHECKLIST}, each item with its reference.}\label{tab:fm:regulators-and-licences:algo}
\end{table}

\section{Tutorial: three offices}\label{tut:fm:regulators-and-licences:tut}

\begin{tutorial}
\textbf{Goal.} Query the map for a proprietary high-frequency firm that starts in London and adds offices in Chicago and Hong Kong.
\textbf{End state:} the regimes each step touches (\cref{tab:fm:regulators-and-licences:profiles}) and the statutory decision periods on the path
(\cref{fig:fm:regulators-and-licences:periods}).
\begin{tutsteps}
\item \textbf{The map.} \texttt{fm\_regmap.rows()} loads the 20 rows of \texttt{data/desk/regmap.csv} (four jurisdictions, ten regimes) and
\texttt{firm.regmap.check} confirms that every row has a reference, an as-of month and a source.
\item \textbf{The profiles.} London: own-account exchange trading with a high-frequency technique. Chicago adds the same in US securities.
Hong Kong adds own-account exchange trading there.
\item \textbf{The queries.} \texttt{fm\_regmap.profiles()} lists the rows, the distinct regimes and the longest statutory decision period.
\item \textbf{The upkeep.} \texttt{firm.regmap.stale(rows, today, 12)} lists rows not checked for a year.
\end{tutsteps}
\end{tutorial}

\begin{table}[ht]
\centering\small
\begin{tabular}{lrrrrl}
\toprule
profile & rows & regimes & required & to check & longest decision \\
\midrule
London & 2 & 1 & 0 & 2 & 6 months (UK) \\
London and Chicago & 5 & 4 & 2 & 2 & 6 months (UK) \\
London, Chicago and Hong Kong & 6 & 5 & 2 & 3 & 6 months (UK) \\
\bottomrule
\end{tabular}
\caption{The regulatory map queried for a firm adding offices (the Chicago step also adds one obligation, the market access rule). Information,
not advice. Data: \texttt{fm\_regmap.profiles}.}\label{tab:fm:regulators-and-licences:profiles}
\end{table}

\begin{omfigure}
\begin{tikzpicture}
\begin{axis}[width=10.5cm,height=4.6cm,xbar stacked,bar width=12pt,xmin=0,xmax=7,xlabel={\small rows of the map touched},ytick={0,1,2},
  yticklabels={London,{London and Chicago},{London, Chicago and Hong Kong}},y dir=reverse,ymin=-0.6,ymax=2.6,
  tick label style={font=\footnotesize},grid=major,grid style={gray!20},table/col sep=comma,
  legend cell align=left,legend style={font=\scriptsize,at={(0.5,-0.35)},anchor=north,/tikz/every even column/.append style={column sep=8pt}},legend columns=3]
\addplot[fill=omProp!55,draw=omProp] table[y=pos,x=required]{figdata/desk/17-regulators-and-licences/profiles.csv};
\addplot[fill=omThm!50,draw=omThm] table[y=pos,x=obligation]{figdata/desk/17-regulators-and-licences/profiles.csv};
\addplot[fill=gray!40,draw=gray] table[y=pos,x=check]{figdata/desk/17-regulators-and-licences/profiles.csv};
\legend{required,obligation,to check}
\end{axis}
\end{tikzpicture}

\omcaption{Rows of the regulatory map a firm touches as it adds offices, by status. Each office adds regimes, and each check row is a question for
the firm's lawyers. Data: \texttt{fm\_regmap.profiles}.}\label{fig:fm:regulators-and-licences:profiles}
\end{omfigure}

\begin{omfigure}
\begin{tikzpicture}
\begin{axis}[width=10.5cm,height=5cm,xbar,bar width=9pt,xmin=0,xmax=400,xlabel={\small days},ytick={0,1,2,3,4,5},
  yticklabels={UK: complete application,UK: incomplete application,EU: complete application,US: grant or open proceedings,US: proceedings concluded,US: with extension},
  y dir=reverse,ymin=-0.6,ymax=5.6,tick label style={font=\footnotesize},grid=major,grid style={gray!20},table/col sep=comma,
  nodes near coords,nodes near coords style={font=\scriptsize}]
\addplot[fill=omDef!55,draw=omDef] table[y=pos,x=days]{figdata/desk/17-regulators-and-licences/periods.csv};
\end{axis}
\end{tikzpicture}

\omcaption{Statutory maximum periods for a regulator's decision on an application, in days (six months counted as 183): UK permission, EU
investment-firm \omterm{def:fm:regulators-and-licences:perimeter}{authorisation}, and US broker-dealer registration.
Data: \texttt{fm\_regmap.PERIODS}; sources in the chapter's ledger.}\label{fig:fm:regulators-and-licences:periods}
\end{omfigure}

The London firm touches one regime, a set of checks on the UK's general prohibition; Chicago adds three (broker-dealer registration, association
membership unless it trades only on its own exchanges, and the market access rule); Hong Kong adds a licence check. The longest statutory decision on
the path is the UK's six months for a complete application, twelve for an incomplete one; the US broker-dealer decision is 45 days, or up to 210 if
the SEC opens proceedings and extends them. The statutory periods start when the application is complete: the months spent writing the programme of
operations, hiring the approved persons and capitalising the entity come first, and the periods are the part of the path a firm does not control.

\begin{remark}[Keeping the map true]\label{rem:fm:regulators-and-licences:upkeep}
Every row of the map is a dated fact about a law that changes. The map's value is in its upkeep: an owner for each jurisdiction, a review date for
each row, and a stale-row report that the compliance function reads. A year after this chapter's check, every row will be stale.
\end{remark}

\section{Build: the regulatory map}\label{bld:fm:regulators-and-licences:regmap}

\begin{build}
\bfield{Purpose} No build in the running-project sense: the subject is law. The analytical tool keeps the map as data, checked and queried.
\bfield{Interface} \texttt{firm.regmap}: \texttt{load}, \texttt{check}, \texttt{query}, \texttt{regimes}, \texttt{stale}, \texttt{longest\_decision},
\texttt{ALGO\_CHECKLIST}; data in \texttt{data/desk/regmap.csv}.
\bfield{Rules} Every row has a reference, an as-of month and a source; statuses are required, obligation or check; nothing is inferred beyond the
rows.
\bfield{Acceptance tests} \texttt{code/firm/regmap/tests/}: queries, regimes and longest periods on a three-row map; the schema check and the stale
report.
\bfield{Stretch} Rows for Singapore, Japan and Australia with their sources; capital requirements per row (chapter 2); a diff between two versions
of the map.
\end{build}

\begin{omsources}
\item Directive 2014/65/EU (MiFID II), articles 2, 4, 5, 7, 17 and 34.
\item Financial Services and Markets Act 2000, sections 19 and 55V.
\item Securities Exchange Act of 1934, section 15; 17 CFR 240.15b9-1; Investment Advisers Act of 1940, section 203; NFA, registration categories.
\item Securities and Futures Commission, licensing: types of regulated activity.
\end{omsources}

\section{Exercises}

\begin{exercise}[$\star$]\label{exo:fm:regulators-and-licences:1}
A firm deals only on its own account in EU equities, through a broker, with no direct electronic access and no algorithm. Is it within the MiFID
II exemption? What changes if it becomes a member of the exchange?
\end{exercise}

\begin{exercise}[$\star$]\label{exo:fm:regulators-and-licences:2}
Name the three features of a high-frequency algorithmic trading technique in MiFID II.
\end{exercise}

\begin{exercise}[$\star$]\label{exo:fm:regulators-and-licences:3}
What is the longest period the SEC can take on a broker-dealer application under section 15(b)(1)?
\end{exercise}

\begin{exercise}[$\star\star$]\label{exo:fm:regulators-and-licences:4}
What does the 2023 amendment of Rule 15b9-1 change for a proprietary firm that trades on an exchange and also off-exchange?
\end{exercise}

\begin{exercise}[$\star\star$]\label{exo:fm:regulators-and-licences:5}
A US firm runs a futures fund for outside investors and gives trading advice. Which CFTC categories does the map point to, and through whom does it
register?
\end{exercise}

\begin{exercise}[$\star\star$]\label{exo:fm:regulators-and-licences:6}
Why does the map mark the UK and Hong Kong own-account rows as checks rather than requirements?
\end{exercise}

\begin{exercise}[$\star\star\star$]\label{exo:fm:regulators-and-licences:7}
\emph{Coding.} Query the map for an EU market maker using a high-frequency technique. How many rows and regimes, and what is the longest statutory
decision?
\end{exercise}

\begin{exercise}[$\star\star\star$]\label{exo:fm:regulators-and-licences:8}
\emph{Find the flaw.} ``We only trade our own money, so we don't need a licence anywhere.''
\end{exercise}

\section{Problem: Three Offices}

\begin{problem}[{Weekend problem --- three offices}]\label{pb:fm:regulators-and-licences:1}
A London proprietary firm plans offices in Chicago and Hong Kong and asks its chief operating officer for the regulatory path.

\medskip\noindent\textbf{Part I --- The perimeter.}
\begin{enumerate}
\item Define the \omterm{def:fm:regulators-and-licences:perimeter}{regulatory perimeter} and \omterm{def:fm:regulators-and-licences:perimeter}{authorisation}.
\item Define \omterm{def:fm:regulators-and-licences:own}{dealing on own account} and an \omterm{def:fm:regulators-and-licences:own}{investment firm}.
\item State the MiFID II own-account exemption and its exceptions.
\item Define \omterm{def:fm:regulators-and-licences:passport}{passporting} and say what it gives an EU-authorised firm.
\end{enumerate}

\medskip\noindent\textbf{Part II --- The United States.}
\begin{enumerate}[resume]
\item Define a \omterm{def:fm:regulators-and-licences:sro}{self-regulatory organisation} and name two.
\item State the broker-dealer registration and membership rules and their timetable.
\item Define a \omterm{def:fm:regulators-and-licences:cftc}{swap dealer}, a \omterm{def:fm:regulators-and-licences:cftc}{commodity pool operator} and an \omterm{def:fm:regulators-and-licences:cftc}{introducing broker}.
\item What does the market access rule add?
\end{enumerate}

\medskip\noindent\textbf{Part III --- The map.}
\begin{enumerate}[resume]
\item What does each row of the map carry, and why?
\item Query the three profiles and give rows and regimes.
\item Give the statutory decision periods on the path.
\item What does the stale-row report show a year later?
\item What does the map not tell the firm?
\end{enumerate}

\medskip\noindent\textbf{Part IV --- The plan.}
\begin{enumerate}[resume]
\item What obligations follow for the firm's algorithms in the EU?
\item What must be done before an application's statutory clock starts?
\item Which step of the path does the firm control, and which not?
\item How would you keep the map up to date?
\item Why is Hong Kong a check and not a requirement in the map?
\item State the \emph{named result}: the regimes the firm's activities touch as it adds a second and a third jurisdiction, and the longest
statutory decision period on its path to trading.
\item In two sentences, write the plan.
\end{enumerate}
\end{problem}

\section{Interview questions}

\begin{interviewq}[$\star$ \iqroles{trader}]\label{iq:fm:regulators-and-licences:1}
Why does a proprietary trading firm that has no clients need a licence in the EU?
\end{interviewq}

\begin{interviewq}[$\star$ \iqroles{developer}]\label{iq:fm:regulators-and-licences:2}
What does it mean, legally, for your firm to engage in algorithmic trading in the EU?
\end{interviewq}

\begin{interviewq}[$\star\star$ \iqroles{risk}]\label{iq:fm:regulators-and-licences:3}
Your firm wants to trade swaps with banks for its own account. What would you check in the US?
\end{interviewq}

\begin{interviewq}[$\star\star$ \iqroles{developer}]\label{iq:fm:regulators-and-licences:4}
Design a data model for a regulatory map that can be audited.
\end{interviewq}

\begin{interviewq}[$\star\star$ \iqroles{trader, risk}]\label{iq:fm:regulators-and-licences:5}
What is \omterm{def:fm:regulators-and-licences:passport}{passporting}, and what happens to it when a country leaves the single market?
\end{interviewq}

\begin{interviewq}[$\star\star\star$ \iqroles{risk}]\label{iq:fm:regulators-and-licences:6}
Plan the regulatory path for a firm moving from trading its own money to managing outside money, in the US and the EU.
\end{interviewq}
